\begin{helper}[Transport from the least active slot]
Let \(G,e_0,k,T\), a generating coloring, and an ETT be fixed.  Let \(q\)
be the last noninitial Lean index, with \(q+1\) equal to the number of
rounds, and assume slot \(q.\mathsf{succ}\) is a PE whose exact chosen vertex
is a terminal support vertex \(v_n\in V(T)\).  Let \(jSlot\) satisfy
\[
 jSlot+1=m(v_n),
\]
where \(m\) is \noderef{ETTFirstAppearanceIndex}.  Put
\(L=\noderef{ExtensionBaseActiveIndices}\).

Assume that whenever \(s,t\in L\), \(s\le t\), their exact chosen vertices
belong to \(V(T)\), and the exact delta field at \(t\) is \(\delta_t\), one
has
\[
 \delta_t\notin\overline{\varphi_s}(T_{s+1}).               \tag{1}
\]
Then every \(t\in L\), with exact chosen vertex in \(V(T)\) and exact delta
field \(\delta_t\), satisfies
\[
 \delta_t\notin
 \overline{\varphi_{jSlot}}(T_{jSlot+1}).                   \tag{2}
\]
The quantified final-tree membership proofs in (1) and (2) are the same
dependent witnesses appearing in the Lean statement.
\end{helper}
\begin{proof}
By \noderef{ExtensionRulesPEFirstAppearanceBound}, the first-appearance
index of the terminal support is at most \(q+1\).  Together with the exact
terminal kind and chosen-vertex equations, this proves directly from
\noderef{ExtensionBaseActiveIndices} that \(q\in L\).  Thus \(L\) is a
nonempty finite set.  Let \(s_0\) be its least Lean index.  For a fixed
\(t\in L\), leastness gives \(s_0\le t\), and hypothesis (1), instantiated
with the exact chosen vertices and delta field, gives
\[
 \delta_t\notin\overline{\varphi_{s_0}}(T_{s_0+1}).         \tag{3}
\]

Membership of \(s_0\) in \(L\), together with
\(jSlot+1=m(v_n)\), implies \(jSlot\le s_0\).  If equality holds, the
coloring and tree indices in (2) and (3) are definitionally the same, and
the result follows.

Assume now that \(jSlot<s_0\), and let \(x\) be any vertex of the displayed
tree at \(jSlot\).  Iterating the ordered-extension fields of
\noderef{TashkinovSeriesStep}, followed by the terminal initial-segment field
of \noderef{ExtendedTashkinovTree}, puts \(x\) in \(V(T)\).  Unfold the
finite minimum in \noderef{ETTFirstAppearanceIndex}.  The ladder entry at
\(jSlot\) is a candidate containing \(x\), so
\[
                         m(x)\le jSlot+1.                    \tag{4}
\]

There is no index \(r\) with \(jSlot\le r<s_0\) whose positive slot is PE
and whose exact chosen vertex is \(x\).  Otherwise (4) and
\(jSlot+1=m(v_n)\) verify the first-appearance clause in the definition of
\(L\); the remaining clauses are the assumed PE kind and exact chosen-vertex
equations.  Hence \(r\in L\), contradicting the leastness of \(s_0\).

Let \(b=s_0-1\).  Apply \noderef{HkuUnsupportedVertexColors} to \(x\) on
the full interval from \(jSlot\) through \(b\).  The preceding paragraph is
exactly its unsupported-vertex premise, and it yields
\[
 \overline{\varphi_{jSlot}}(x)
   =\overline{\varphi_{s_0}}(x).                             \tag{5}
\]
If \(\delta_t\) were tree-missing at the displayed tree for \(jSlot\),
\noderef{TreeMissingColors} would supply such a vertex \(x\) missing it.
Equation (5) preserves that missing color, while the iterated ordered
initial segments put \(x\) in the displayed tree for \(s_0\).  Thus
\(\delta_t\) would be tree-missing there, contradicting (3).  This proves
(2).
\end{proof}
